\section{Characterization of almost sure convergence}
\label{sec: borel cantelli}


In Theorem \ref{thm: convergence implications} we have
seen that almost sure convergence implies convergence in probability. In this
section we will discover that the speed at which the term
\begin{equation*}
    \prob*{\norm{X_n - X} \ge \epsilon}
\end{equation*}
converges to zero determines whether we also have almost sure convergence.
Roughly speaking if this probability decays like \(\frac1n\) or slower, then
we do not have almost sure convergence. But if it decays faster, then we do.
The key results to formalize this claim are known as the Borel-Cantelli lemmas
(Theorem \ref{thm: borel-cantelli}). With these we will characterize almost sure
convergence in terms of speed of convergence in Corollary \ref{cor:
characterization of almost sure convergence}.

To state the Borel-Cantelli lemmas we recall some definitions about limits of
events.

\begin{definition}[Limits of events]
    If the sequence of events \((A_n)_{n=1}^\infty\) 
    \begin{itemize}
        \item is \emph{increasing}, that is \(A_1\subseteq A_2\subseteq \ldots\),
        then we write
        \[
            \lim_{n\to \infty} A_n \coloneq \bigcup_{n\in \nat} A_n
            \qquad \text{ or } \quad A_n \uparrow \bigcup_{n\in \nat} A_n.
        \]
        \item is \emph{decreasing}, that is \(A_1\supseteq A_2\supseteq \ldots\),
        then we write
        \[
            \lim_{n\to \infty} A_n \coloneq \bigcap_{n\in \nat} A_n
            \qquad \text{ or } \quad A_n \downarrow \bigcap_{n\in \nat} A_n.
        \]
    \end{itemize}
\end{definition}

\begin{remark}[Continuity of the probability measure]
    \label{rem: continuity of probability}
    Recall that the probability measure is continuous for increasing and decreasing
    sequences of events \continuityOfMeasure, that is
    \begin{equation}
        \label{eq: continuity of probability}
        \prob[\big]{\lim_{n\to \infty} A_n}
        = \lim_{n\to \infty} \prob{A_n}.
    \end{equation}
    We have already used this in the proof that almost sure convergence implies
    convergence in probability in Theorem \ref{thm: convergence implications}.
\end{remark}

Now recall that an event \(A\) ``happens'' if \(\omega \in A\). This
motivates the following definition.

\begin{definition}[``Infinitely many'' and ``eventually all'']
    Let \((A_n)_{n=1}^\infty\) be a sequence of events. Then
    \begin{itemize}
        \item 
        the \emph{limes superior}
        is the event that infinitely many of the events \(A_n\) happen
        \[\begin{aligned}
            \limsup_{n\to\infty} A_n 
            &\coloneq \bigcap_{N\in \nat} \bigcup_{n \ge N} A_n,
            \\
            &= \set{\omega \in \Omega : \omega \in A_n \text{ for infinitely many } n}.
        \end{aligned}
        \]
        
        \item
        the \emph{limes inferior}
        is the event that eventually all of the events \(A_n\) happen
        \[\begin{aligned}
            \liminf_{n\to\infty} A_n 
            &\coloneq \bigcup_{N\in \nat} \bigcap_{n \ge N} A_n,
            \\
            &= \set{\omega \in \Omega : \omega \in A_n \text{ for all \(n\) after some } N\in \nat}.
        \end{aligned}
        \]
    \end{itemize}
\end{definition}

To better understand these definitions it is helpful to convince yourself of the
following properties and compare the limes superior/inferior for sets with
its analogue for real-valued sequences.

\begin{lemma}[Liminf and limsup properties]
    \label{lem: liminf and limsup properties}
    We have
    \begin{enumerate}[label=(\roman*)]
        \item \(\liminf_n A_n \subseteq \limsup_n A_n\),
        \item \((\limsup_n A_n)^\complement = \liminf_n A_n^\complement\),
        \item \(\limsup_n \ind{A_n}(\omega) = \ind*{\limsup_n A_n}(\omega)\) for all \(\omega \in \Omega\),
        \item \(\liminf_n \ind{A_n}(\omega) = \ind*{\liminf_n A_n}(\omega)\) for all \(\omega \in \Omega\).
    \end{enumerate}
\end{lemma}
\begin{proof}
    Exercise \ref{ex: liminf and limsup properties}. Note that indicators are always in \(\{0,1\}\)
    and use the definition of the limes superior and inferior.
\end{proof}

\begin{mdframed}[innertopmargin=0pt]
\begin{theorem}[Borel-Cantelli Lemmas]
    \label{thm: borel-cantelli}
    Let \((A_n)_{n=1}^\infty\) be a sequence of events. Then
    \begin{alignat}{3}
        \tag{BC-1}\label{eq: borel cantelli 1}
        \sum_{n=1}^\infty \prob*{A_n}<\infty& \quad & \implies & \quad & \prob[\Big]{\limsup_{n\to \infty} A_n}=0.
    \intertext{
        If the events \(A_n\) are further \textbf{independent} events, then
    }
        \tag{BC-2}\label{eq: borel cantelli 2}
        \sum_{n=1}^\infty \prob{A_n}=\infty& &\implies& &\prob[\Big]{\limsup_{n\to \infty} A_n}=1.
    \end{alignat}
\end{theorem}
\end{mdframed}


\begin{proof}[Proof of Borel-Cantelli (Theorem \ref{thm: borel-cantelli})]
Let \((A_n)_{n=1}^\infty\) be a sequence of events with
\(\sum_{n=1}^\infty \prob*{A_n} < \infty\). Observe that
a smaller and smaller union of events is a decreasing sequence and consequently
\begin{align*}
    \prob[\Big]{\limsup_{n\to \infty} A_n}
    = \prob[\Big]{ \lim_{N\to\infty} \bigcup_{n \ge N} A_n }
    &= \lim_{N\to\infty} \prob[\Big]{\bigcup_{n \ge N} A_n}
    \\
    &\le \lim_{N\to\infty} \sum_{n=N}^\infty \prob{A_n}
    = 0.
\end{align*}
This proves the first Borel-Cantelli lemma \eqref{eq: borel cantelli 1}.
For the second Borel-Cantelli lemma \eqref{eq: borel cantelli 2}
let $(A_n)_{n=1}^\infty$ be a sequence of independent events. By
Lemma \ref{lem: liminf and limsup properties} we have
\[\begin{aligned}
    \prob[\Big]{\bigl(\limsup_{n\to\infty}A_n\bigr)^\complement}
    &= \prob[\Big]{\bigcup_{N\in \nat} \bigcap_{n \ge N} A_n^\complement}
    \\
    \overset{\eqref{eq: continuity of probability}}&= \lim_{N\to\infty} \prob[\Big]{\bigcap_{n \ge N} A_n^\complement}
    \\
    \overset{\eqref{eq: continuity of probability}}&= \lim_{N\to\infty} \lim_{k\to \infty} \prob[\Big]{\bigcap_{n=N}^{k} A_n^\complement}
    \\
    \overset{\text{indep.}}&= \lim_{N\to\infty} \lim_{k\to \infty} \prod_{n=N}^{k} (1-\prob{A_n})
    \\
    \overset{1+x\le e^x}&\le \lim_{N\to\infty} \underbrace{\lim_{k\to \infty} \exp\Bigl(-\sum_{n=N}^{k} \prob{A_n}\Bigr)}_{=0} = 0.
\end{aligned}
\]
Where we used that \(\sum_{n=1}^\infty \prob{A_n} = \infty\) and \(\exp(-x)\to 0\) for \(x\to \infty\).
\end{proof}

\begin{remark}[Pairwise independence]
    \label{rem: pairwise independence}
    With more effort it is possible to generalize the
    Borel-Cantelli Lemma 2 \eqref{eq: borel cantelli 2} to require only
    \underline{pairwise} independence \citep{doeringStochastik2022}.
\end{remark}

The Borel-Cantelli lemmas now lead to the following characterization of almost
sure convergence.

\begin{mdframed}[innertopmargin=0pt]
\begin{corollary}[Characterization of almost sure convergence]
    \label{cor: characterization of almost sure convergence}
    Let \((X_n)_{n=1}^\infty\) be a sequence of random variables and \(X\) another random variable.
    Then
    \[
        \forall \epsilon > 0: \quad
        \sum_{n=1}^\infty \prob[\big]{\norm{X_n - X} \ge \epsilon} < \infty
         \quad \implies \quad \underbrace{
            \prob[\Big]{\lim_{n\to\infty} X_n = X} =1
        }_{
            \iff X_n\overset{\as}\to X
        }.
    \]
    If the random variables \((X_n-X)_{n\in \nat}\) are (pairwise)\footnote{see Remark \ref{rem: pairwise independence}} independent, then
    \[
        \exists \epsilon > 0: \quad
        \sum_{n=1}^\infty \prob[\big]{\norm{X_n - X} \ge \epsilon} = \infty
        \quad \implies \quad \prob[\Big]{\lim_{n\to\infty} X_n = X} =0.
    \]
    In the independent case this implies a \(0\)-\(1\) law: If the probability of
    convergence is greater zero it must be one!
\end{corollary}
\end{mdframed}

This corollary shows that almost sure convergence simply means that we have
``fast enough'' convergence in probability such that sequence of
\(\prob[\big]{\norm{X_n - X} \ge \epsilon}\) does not only converge to zero but is also
summable.

\begin{proof}
    We have by definition of convergence and continuity of measure (Remark \ref{rem: continuity of probability})
    \begin{equation}
    \label{eq: characterization of does not converge}    
    \begin{aligned}[b]
        \prob[\Big]{\lim_{n\to\infty} X_n \neq X}
        &= \prob[\Big]{ \bigcup_{m \in \nat} \bigcap_{N\in \nat}\bigcup_{n \ge N} \set{\norm{X_n - X} \ge \tfrac1m} }
        \\
        &= \lim_{m\to \infty} \prob[\Big]{ \limsup_{n\to \infty} \{\norm{X_n - X} \ge \epsilon\} }.
    \end{aligned}
    \end{equation}
    Now if for all \(\epsilon > 0\) we have that \(\sum_{n=1}^\infty \prob[\big]{\norm{X_n - X} \ge \epsilon} < \infty\),
    then by the first Borel-Cantelli lemma \eqref{eq: borel cantelli 1}
    we have that the probability in \eqref{eq: characterization of does not converge} is zero
    and consequently we get almost sure convergence.

    Now if the random variables \((X_n - X)_{n\in \nat}\) are independent, then
    the sets \(\set{\norm{X_n - X} \ge \epsilon}\) are independent events. If there is an
    \(\epsilon > 0\) such that \(\sum_{n=1}^\infty \prob[\big]{\norm{X_n - X} \ge \epsilon} = \infty\),
    then the application of the second Borel-Cantelli lemma \eqref{eq: borel
    cantelli 2} to this particular \(\epsilon\) shows that the probability
    in \eqref{eq: characterization of does not converge} must be one.
\end{proof}

An application of this characterization of almost sure convergence is the
insight that from every sequence that converges in probability we can extract
a subsequence that converges almost surely.

\begin{lemma}[Subsequences converge almost surely]
    \label{lem: subsequences converge almost surely}
    Let \((X_n)_{n\in \nat}\) be a sequence of random variables and \(X\)
    another random variable such that \(X_n \overset{p}\to X\).
    Then there exists a subsequence \((X_{n_k})_{k\in \nat}\)
    such that \(X_{n_k} \overset{\as}\to X\).
\end{lemma}

The idea of the proof is that we can increase the speed of convergence
of \(\prob[\big]{\norm{X_n - X} \ge \epsilon}\) by choosing a subsequence.
For example in the case
\[
    \prob[\big]{\norm{X_n - X} \ge \epsilon} \le \frac{C_\epsilon}{n},
\]
we can choose \(n_k = k^2\) such that for every \(\epsilon > 0\)
\begin{equation}
    \label{eq: subsequence summable}    
    \sum_{k=1}^\infty \prob[\big]{\norm{X_{n_k} - X} \ge \epsilon}
    \le \sum_{k=1}^\infty \frac{C_\epsilon}{k^2} < \infty.
\end{equation}
Corollary \ref{cor: characterization of almost sure convergence} then
implies that \(X_{n_k} \overset{\as}\to X\).
The actual proof is a little bit more complicated for two reasons: First, we do
not have an explicit rate of convergence. And second, we do not want to select a
subsequence that depends on \(\epsilon\) since it must work for all \(\epsilon >
0\) simultaneously.


\begin{proof}
    By convergence in probability we have for every \(\epsilon >0\),
    in particular \(\epsilon =\frac1k\), that
    there exists \(N_k\) such that for all \(n \ge N_k\)
    \[
        \prob{\norm{X_n - X} \ge \tfrac1k} \le \tfrac1{2^k}.
    \]
    By choosing the subsequence \(n_k = N_k\) we thus get
    for every \(\epsilon > 0\)
    \[
        \sum_{k=1}^\infty \prob[\big]{\norm{X_{n_k} - X} \ge \epsilon}
        = \underbrace{\sum_{k=1}^{\lceil \frac1\epsilon \rceil-1}
        \prob[\big]{\norm{X_{n_k} - X} \ge \epsilon}}_{< \infty}
        +\sum_{k=\lceil \frac1\epsilon \rceil}^\infty
        \prob[\big]{\norm{X_{n_k} - X} \ge \epsilon}
    \]
    For the remaining summands we have \(\epsilon > \frac1k\) and therefore
    \[
       \sum_{k=\lceil \frac1\epsilon \rceil}^\infty
        \prob[\big]{\norm{X_{n_k} - X} \ge \epsilon}
        \le \sum_{k=\lceil \frac1\epsilon \rceil}^\infty
        \prob[\big]{\norm{X_{n_k} - X} \ge \tfrac1k}
        \le \sum_{k=1}^\infty \frac1{2^k} = 1 < \infty.
    \]
    Consequently, by Corollary \ref{cor: characterization of almost sure convergence}
    we have that \(X_{n_k} \overset{\as}\to X\).
\end{proof}

\subsection*{Exercises \thesection}


\begin{exercise}[Liminf and Limsup Properties \Coffeecup]
    \label{ex: liminf and limsup properties}
    Prove Lemma \ref{lem: liminf and limsup properties}. 

    \begin{solution}
        \begin{enumerate}[label=(\roman*)]
            \item To show that  \(\liminf_n A_n \subseteq \limsup_n A_n\) let
            \[
                \omega \in \liminf_n A_n = \bigcup_{N\in \nat} \bigcap_{n\ge N} A_n.
            \]
            then there exists an \(N\) such that for all \(n \ge N\) we have
            \(\omega \in A_n\). In particular, there are infinitely many \(n\)
            such that \(\omega \in A_n\) and consequently for all \(N\in \nat\) there
            exists \(n\ge N\) such that \(\omega \in A_n\), which implies
            \[
                \omega \in \bigcap_{N\in \nat} \bigcup_{n\ge N} A_n = \limsup_n A_n.
            \]
            
            \item\label{it: de morgans on liminf/limsup} We have using De Morgan's laws that
            \[
            \begin{aligned}
                (\limsup_n A_n)^\complement
                &= \Bigl(\bigcap_{N\in \nat} \bigcup_{n\ge N} A_n\Bigr)^\complement
                = \bigcup_{N\in \nat} \Bigl(\bigcup_{n\ge N} A_n\Bigr)^\complement
                = \bigcup_{N\in\nat} \bigcap_{n\ge N} A_n^\complement
                \\
                &= \liminf_n A_n^\complement.
            \end{aligned}
            \]

            \item\label{it: function limsup to set limsup} To prove \(\limsup_n \ind{A_n} = \ind*{\limsup_n A_n}\) 
            observe that 
            \[
                \limsup_n \ind{A_n}(\omega)
                = \lim_{N\to \infty} \sup_{n\ge N} \ind{A_n}(\omega)
            \]
            is in the set \(\set{0,1}\) for all \(\omega\) since
            \(\ind{A_n}(\omega) \in \set{0,1}\) for all \(n\) and
            the supremum is thereby either \(0\) or \(1\). We therefore only need to show
            that \(\limsup_n \ind{A_n}(\omega) = 1\) if and only if
            \begin{equation}
                \label{eq: omega in limsup_n An} 
                \omega \in \limsup_n A_n = \bigcap_{N\in \nat} \bigcup_{n\ge N} A_n.
            \end{equation}
            Assume \eqref{eq: omega in limsup_n An} holds, then for all \(N\in
            \nat\) there exists an \(n \ge N\) such that \(\omega \in A_n\). In particular,
            since for any \(N\in \nat\) we can find \(n\ge N\) such that \(\ind{A_n}(\omega) = 1\)
            and therefore \(\sup_{n\ge N} \ind{A_n}(\omega) = 1\). Conversely, if
            \(\limsup_n \ind{A_n}(\omega) = 1\), then for all \(N\in \nat\) we have
            \(\sup_{n\ge N} \ind{A_n}(\omega) =1\) and consequently there must exist
            \(n\ge N\) such that \(\ind{A_n}(\omega) = 1\) and therefore
            \(\omega \in A_n\). This implies \eqref{eq: omega in limsup_n An}.

            \item Observe that we can translate this to the previous results as follows:
            \[\begin{aligned}[b]
                \liminf_n \ind{A_n}
                &= \liminf_n (1-\ind{A_n^\complement})
                \\
                &= 1 - \limsup_n \ind{A_n^\complement}
                \\
                \overset{\text{\ref{it: function limsup to set limsup}}}&= 1 - \ind*{\limsup_n A_n^\complement}
                \overset{\text{\ref{it: de morgans on liminf/limsup}}}= \ind*{\liminf_n A_n}.
            \end{aligned}
            \qedhere
            \]
        \end{enumerate}        
    \end{solution}
\end{exercise}


\begin{exercise}[Borel's \(0\)-\(1\) law \Coffeecup]
    \label{ex: borel cantelli 0-1 law}
    Let \((A_n)_{n=1}^\infty\) be a sequence of independent events. Prove that 
    the probability of their liminf and limsup are either zero or one, that is
    \[
        \prob[\Big]{\limsup_{n\to \infty} A_n} \in \{0,1\}
        \quad\text{and}\quad
        \prob[\Big]{\liminf_{n\to \infty} A_n} \in \{0,1\}.
    \]
    Consequently, if the probability is greater zero it must be one!

    \noindent\textbf{Hint:} Use the Borel-Cantelli Lemmas (Theorem \ref{thm: borel-cantelli})
    and Lemma \ref{lem: liminf and limsup properties}. Without justification you
    may use the measure theoretic fact that \((A_n^\complement)_{n\in \nat}\)
    is indpendent if \((A_n)_{n\in \nat}\) is independent.
    \begin{remark}[Kolmogorov's \(0\)-\(1\) law]
        \label{rem: kolmogorov 0-1 law}
        The Borel-Cantelli \(0\)-\(1\) law is a special case of Kolmogorov's \(0\)-\(1\) law
        about ``tail events'' \citep[Thm.~2.37, Cor.~2.38]{klenkeProbabilityTheoryComprehensive2014},
        as the liminf and limsup are such ``tail events''.
    \end{remark}
\end{exercise}

\begin{solution}
    Observe that there are only two possibilities, either
    \[
        \sum_{n=1}^\infty \prob{A_n} < \infty
            \quad\text{or}\quad
        \sum_{n=1}^\infty \prob{A_n} = \infty.
    \]
    In the first case we have by the first Borel-Cantelli lemma \eqref{eq: borel cantelli 1} that
    \[
        \prob[\Big]{\limsup_{n\to \infty} A_n} = 0.
    \]
    in the second case we have by the second Borel-Cantelli lemma \eqref{eq: borel cantelli 2} that
    \[
        \prob[\Big]{\limsup_{n\to \infty} A_n} = 1.
    \]
    Consequently, we have the claim for \(\limsup_n A_n\). Due to
    \[
        \liminf_n A_n = \limsup_n A_n^\complement
    \]
    by Lemma \ref{lem: liminf and limsup properties} and the independence of
    \(A_n^\complement\), the
    same arguments apply to \(A_n^\complement\) so we also have the claim for
    \(\liminf_n A_n\).
    \qedhere
\end{solution}


\begin{exercise}[subtitle={Almost sure but not \(L^p\) convergence \Coffeecup}]
    \label{ex: as but not Lp}
    Let \(U\sim \uniform(0,1)\) and define
    \[
        X_n = 2^n \ind{[0, 2^{-n})}(U).
    \]
    Show that \(X_n \overset{\as}\to 0\) but \(X_n\) does not converge to \(0\)
    in \(L^1\) and hence in no other \(L^p\) space with \(p \ge 1\).
\end{exercise}

\begin{solution}
    First, we show that \(X_n \overset{\as}\to 0\). For any \(\epsilon > 0\)
    we have
    \[\begin{aligned}
        \sum_{n=0}^\infty \prob{\abs{X_n - 0} \ge \epsilon}
        &\le \sum_{n=0}^\infty \prob{X_n \ge \min\set{\epsilon, 1}}
        \\
        &= \sum_{n=0}^\infty\prob{U < 2^{-n}} = \sum_{n=0}^\infty 2^{-n} < \infty.
    \end{aligned}
    \]
    Consequently, by Corollary \ref{cor: characterization of almost sure convergence} we have \(X_n \overset{\as}\to 0\).
    To see that \(X_n\) does not converge to \(0\) in \(L^1\) we observe
    \[
        \E{\abs{X_n - 0}}
        = \E{X_n}
        = 2^n \prob{U < 2^{-n}} = 1
        \not\to 0.
        \qedhere
    \]
\end{solution}


\begin{exercise}[Infinitely often or never successful \Coffeecup]
    Let \(X_n \sim \bernoulli(n^{-r})\) be independent Bernoulli random variables,
    whose success probability decreases in \(n\) at rate \(r > 0\). Note that we say
    that \(X_n\) is a success if \(X_n = 1\). Show that
    \begin{enumerate}[label=(\alph*)]
        \item for \(r > 1\) there are finitely many successes with probability one.
        That is after some random but finite time \(N\in \nat\) all \(X_n\) are zero.
        \begin{solution}
            Since we have
            \[
                \sum_{n=1}^\infty \prob{X_n = 1} = \sum_{n=1}^\infty n^{-r} < \infty,
            \]
            we have by the first Borel-Cantelli lemma \eqref{eq: borel cantelli 1} that
            \[
                \prob{X_n=1 \text{ for infinitely many } n} = \prob{\limsup_{n\to \infty} \set{X_n = 1}} = 0.
                \qedhere
            \]
        \end{solution}

        \item for \(r \le 1\) there are infinitely many successes with probability one.
        \begin{solution}
            Since we have
            \[
                \sum_{n=1}^\infty \prob{X_n = 1} = \sum_{n=1}^\infty n^{-r} = \infty,
            \]
            we have by the second Borel-Cantelli lemma \eqref{eq: borel cantelli 2} that
            \[
                \prob{X_n=1 \text{ for infinitely many } n} = \prob{\limsup_{n\to \infty} \set{X_n = 1}} = 1.
                \qedhere
            \]
        \end{solution}
    \end{enumerate}
\end{exercise}


\begin{exercise}[Further characterize almost sure convergence \Coffeecup]
    Prove the following are equivalent
    \begin{enumerate}[label=(\roman*), noitemsep]
        \item\label{it: as convergence} \(X_n \overset{\as}\to X\), that is \(\prob{X_n \to X} =1\)
        \item\label{it: limsup set convergence} \(\prob[\big]{\limsup_{n \to \infty} \set{|X_n - X| \ge \epsilon}} = 0\) for all \(\epsilon > 0\)
        \item\label{it: limsup convergence} \(\prob[\big]{\limsup_{n \to \infty} |X_n - X| \ge \epsilon} = 0\) for all \(\epsilon >0\)
        \item\label{it: sup convergence} \(\lim_{N\to \infty} \prob[\big]{\sup_{n \ge N} |X_n - X| \ge \epsilon} = 0\) for all \(\epsilon > 0\).
    \end{enumerate}
\end{exercise}

\begin{solution}
    First we prove \ref{it: as convergence} \(\Rightarrow\) \ref{it: limsup set convergence}.
    By definition of convergence we have
    \[\begin{aligned}
        0 = \prob{X_n \not\to X}
        &= \prob[\Big]{\bigcup_{\epsilon >0} \bigcap_{N\in \nat} \bigcup_{n\ge N} \set{|X_n - X| \ge \epsilon}}
        \\
        \overset{\text{fix } \epsilon}&\ge \prob[\Big]{\bigcap_{N\in \nat} \bigcup_{n\ge N} \set{|X_n - X| \ge \epsilon}}
        \\
        \overset{\text{def.}}&= \prob[\big]{\limsup_{n \to \infty} \set{|X_n - X| \ge \epsilon}}.
    \end{aligned}
    \]
    For the converse direction \ref{it: limsup set convergence} \(\Rightarrow\) \ref{it: as convergence}, we use
    \[\begin{aligned}
        \prob{X_n \not\to X}
        &= \prob[\Big]{\bigcup_{m\in \nat} \bigcap_{N\in \nat} \bigcup_{n\ge N} \set{|X_n - X| < \tfrac1m}}
        \\
        &= \lim_{m\to \infty} \prob[\Big]{\limsup_{n \to \infty} \set{|X_n - X| \ge \tfrac1m} = 0} = 0.
    \end{aligned}
    \]
    The last equation follows from the assumption as \(\epsilon=\frac1m\) are just special cases.
    Next we show \ref{it: limsup set convergence} \(\Leftrightarrow\) \ref{it: limsup convergence}.
    Observe that by definition of the limsup of sets
    \[\begin{aligned}
        \prob[\big]{\limsup_{n \to \infty} \set{\abs{X_n - X} \ge \epsilon}}
        &= \prob[\big]{\forall N \in \nat,\; \exists n\ge N: \abs{X_n - X} \ge \epsilon}.
        \\
        &= \prob[\Big]{\forall N \in \nat: \sup_{n\ge N} \abs{X_n - X} \ge \epsilon}.
        \\
        &= \prob[\big]{\limsup_{n \to \infty} |X_n - X| \ge \epsilon}.
    \end{aligned}
    \]
    Consequently the two statements are equivalent. Next we show
    \ref{it: limsup convergence} \(\Leftrightarrow\) \ref{it: sup
    convergence}.  For this observe that the sets \(\set{\sup_{n \ge N} |X_n
    - X| \ge \epsilon}\) are a decreasing sequence in \(N\) and therefore
    \begin{align*}
        \prob[\big]{\limsup_{n \to \infty} |X_n - X| \ge \epsilon}
        &= \prob[\Big]{\bigcap_{N\in \nat} \set{\sup_{n \ge N} |X_n - X| \ge \epsilon}}
        \\
        &= \lim_{N\to \infty} \prob[\big]{\sup_{n \ge N} |X_n - X| \ge \epsilon}.
        \qedhere
    \end{align*}
\end{solution}


\begin{exercise}[Infinite Monkey Theorem \Coffeecup\Coffeecup]
    The infinite monkey theorem states that an immortal monkey hitting keys at
    random on a typewriter keyboard for an infinite amount of time will
    eventually type any arbitrarily long text -- for example all the
    works of William Shakespeare -- with probability one.  It is also known as
    the ``law of truly large numbers''.
    Prove it!
    \begin{enumerate}[label={\emph{Level \arabic*:}},ref={Level \arabic*},leftmargin=4em,noitemsep]
        \item Assume the hit keys are independent and uniformly distributed.
        \item\label{it: level iid} Assume the hit keys are iid with a probability distribution that
        assigns each of the (finite) keys a positive probability.
        \item\label{it: level markov chain} Assume the hit keys are a stationary Markov chain whose transition
        probabilities from key \(x\) to key \(y\) are positive for all pairs
        \(x,y\) from the (finite) alphabet.
    \end{enumerate}
    \textbf{Hint:} \ref{it: level markov chain} is significantly harder than
    the previous levels. For it you already need to know Markov chains
    and modify the proof of the second Borel-Cantelli lemma for your purposes.
    You may want to come back to this later.
\end{exercise}

\begin{solution}
    Let \(\alphabet\) be a finite alphabet of keys and let \(X_n\in
    \alphabet\) be the letter typed at time \(n\). Let \(y\in \alphabet^m\)
    be the desired text of length \(m\). We essentially assume \(m=1\)
    without loss of generality with the definition of
    \[
        Y_n = (X_{nm}, X_{nm+1}, \ldots, X_{m(n+1)-1}) \in \alphabet^m,
    \]
    which is still a sequence of iid random variables in \ref{it: level iid}
    (or still a Markov chain in \ref{it: level markov chain}).  Let us consider \ref{it: level iid}
    first. Since
    \[\begin{aligned}
        \delta \coloneq \prob{Y_n = y}
        &= \prob{X_{nm} = y_1, X_{nm+1} = y_2, \ldots, X_{m(n+1)-1} = y_m}
        \\
        &= \prod_{i=1}^m \prob{X_{nm+i-1} = y_i} > 0,
    \end{aligned}
    \]
    Due to \(\sum_{n=1}^\infty \prob{Y_n = y} = \sum_{n=1}^\infty \delta = \infty\)
    we may apply the second Borel-Cantelli lemma \eqref{eq: borel cantelli 2} to get
    \[
        \prob[\big]{Y_n = y \text{ for infinitely many } n}
        = \prob[\big]{\limsup_{n\to\infty} \set{Y_n = y}}
        \overset{\text{\eqref{eq: borel cantelli 2}}}= 1.
    \]
    Observe that we have proven more than the desired claim: The monkey will
    not only type the desired text \(y\) at least once but infinitely many
    times with probability one!

    For the proof of \ref{it: level markov chain} we need to avoid independence.
    Since any transition probability is positive, we have
    \[
        \delta \coloneq \prob{Y_n = y \mid Y_{n-1} \neq y}
        = \sum_{\substack{x \in \alphabet^m \\ x \neq y}}\prob{Y_n = y \mid Y_{n-1} = x}
        > 0.
    \]
    We will now modify the proof of the second Borel-Cantelli lemma 
    in the place where independece was used. Specifically, we have
    \begin{equation}
        \label{eq: modified borel cantelli 2}
        \prob[\big]{(\limsup_{n\to\infty} \set{Y_n = y})^\complement}
        = \lim_{N\to \infty} \lim_{k\to \infty} \prob[\Big]{\bigcap_{n=N}^k \set{Y_n \neq y}}.
    \end{equation}
    with
    \[\begin{aligned}
        \prob[\Big]{\bigcap_{n=N}^k \set{Y_n \neq y}}
        &= \prod_{n=N}^k \prob{Y_n \neq y \mid Y_{n-1} \neq y,\ldots, Y_{N} \neq y}
        \\
        \overset{\text{Markov}}&= \prob{Y_n\neq y}\prod_{n=N+1}^k \prob{Y_n \neq y \mid Y_{n-1} \neq y}
        \\
        &\le (1-\delta)^{k-N-1}.
    \end{aligned}
    \]
    Thereby \eqref{eq: modified borel cantelli 2} is zero. Consequently,
    the probability that \(Y_n = y\) happens infinitely often is one.
\end{solution}

